\section{Forschungsstand und theoretische Grundlagen}
\label{sec:forschungsstand-und-theoretische-grundlagen}

In sozialen Gruppen existieren Meinungsführende (Opinion-Leaders), die als Fachleute sowie Ratgebende für bestimmte Themen gelten und die im Rahmen interpersonaler Kommunikation ihre Bewertungen, Interpretationen oder Meinungen an die Meinungsfolgenden (Opinion-Followers) weitergeben (Lazarsfeld, Berelson \& Gaudet, 1948). Daneben bezeichnen ‚virtuelle Meinungsführer‘ Medienpersönlichkeiten, die hohes Vertrauen und Ansehen genießen, mit denen jedoch nicht interpersonal kommuniziert wird (Katz, 1957).

Infolge der erhöhten Nutzung sozialer Medien wird die Gültigkeit der Theorie jedoch infrage gestellt: Bennett und Manheim (2006) stellten fest, dass die Grenzen zwischen den Medien, den Meinungsführenden und der Öffentlichkeit bei sozialen Medien verschwimmen. Thorson und Wells (2015) entwarfen daher ein neues Modell, das explizit den Informationsfluss in sozialen Netzwerken beschreibt.

Wie sich der Two-Step-Flow und das Konzept des Opinion-Leadership auf Twitter anwenden lassen, wurde noch nicht in der Literatur behandelt und der Aspekt der politischen Kommunikation wurde ebenfalls weitgehend vernachlässigt.